\chapter{Ekspandery: ekspansja, widmo Laplasjanu i spacery losowe}
\label{ch:tg-expanders}

Rozdział \ref{ch:tg-random-graphs} pokazał, jak własności
grafów badać probabilistycznie. Teraz zajmiemy się grafami,
w których mały zbiór wierzchołków ma wiele połączeń
z dopełnieniem. Powiążemy tę cechę z widmem macierzy
i szybkim mieszaniem spacerów losowych.
Wszystkie oznaczenia i założenia potrzebne w twierdzeniach
podajemy jawnie, w szczególności rozróżniamy
ekspansję krawędziową i wierzchołkową.

\section{Przekroje i stałe ekspansji}
\label{sec:tg-expanders-expansion}

Niech $G=(V,E)$ będzie skończonym grafem prostym
nieskierowanym, o $n=|V|\geq2$ wierzchołkach.
Przez $\partial_E S$ oznaczamy zbiór krawędzi
mających dokładnie jeden koniec w $S\subseteq V$,
zaś przez $\partial_V S=N_G(S)\setminus S$
zbiór sąsiadów zbioru $S$ leżących poza nim.

\begin{definicja}[Ekspansja krawędziowa i wierzchołkowa]
\label{def:tg-expander-constants}
Niech $G=(V,E)$ będzie skończonym grafem prostym
nieskierowanym, $n=|V|\geq2$. Definiujemy
\[
h_E(G)=
\min_{\substack{\varnothing\ne S\subseteq V\\|S|\leq n/2}}
\frac{|\partial_E S|}{|S|},
\qquad
h_V(G)=
\min_{\substack{\varnothing\ne S\subseteq V\\|S|\leq n/2}}
\frac{|\partial_V S|}{|S|}.
\]
Jeżeli ponadto $G$ jest $d$-regularny, gdzie $d\geq1$,
to jego \emph{stałą ekspansji znormalizowanej}
(izoperymetryczną) nazywamy
$\Phi(G)=h_E(G)/d$.
\end{definicja}

\begin{stwierdzenie}[Porównanie dwóch rodzajów ekspansji]
\label{prop:tg-expander-edge-vertex}
Jeżeli $G$ jest skończonym grafem prostym
nieskierowanym o co najmniej dwóch wierzchołkach
i maksymalnym stopniu $\Delta\geq1$, to
\[
h_V(G)\leq h_E(G)\leq\Delta\,h_V(G).
\]
W szczególności dla grafu $d$-regularnego
$\Phi(G)\leq h_V(G)\leq d\,\Phi(G)$.
\end{stwierdzenie}
\begin{proof}
Dla każdego dopuszczalnego $S$
każdy wierzchołek $\partial_VS$
jest końcem co najmniej jednej
i co najwyżej $\Delta$ krawędzi $\partial_ES$.
Zatem $|\partial_VS|\leq|\partial_ES|
\leq\Delta|\partial_VS|$.
Biorąc minimum pierwszej nierówności
po wszystkich $S$, otrzymujemy $h_V\leq h_E$.
Dla drugiej bierzemy zbiór $S$ realizujący
minimum definiujące $h_V$ i dostajemy
$h_E\leq|\partial_ES|/|S|
\leq\Delta h_V$.
Dzieląc pierwszą i drugą nierówność przez $d$
dla $\Delta=d$, otrzymujemy
$h_V/d\leq\Phi\leq h_V$.
Zapisując je równoważnie, uzyskujemy
$\Phi\leq h_V\leq d\Phi$.
\end{proof}

\begin{uwaga}[Wymóg spójności]
W grafie niespójnym składowa o najmniejszej liczności
ma co najwyżej $n/2$ wierzchołków.
Jej granica krawędziowa jest pusta;
dlatego $\Phi(G)=0$ dla każdego niespójnego
grafu $d$-regularnego. Odwrotnie, jeśli graf jest spójny, to każdy
niepusty właściwy podzbiór ma niepustą granicę.
Minimum po skończenie wielu dopuszczalnych zbiorach
jest więc dodatnie. Dla $d\geq1$ regularność daje
równoważność spójności i warunku $\Phi(G)>0$.
\end{uwaga}

\begin{definicja}[Rodzina ekspanderów regularnych]
\label{def:tg-expander-family}
Rodzinę grafów prostych $(G_j)_{j\geq1}$
nazywamy \emph{rodziną ekspanderów
o ograniczonym stopniu}, jeśli
$|V(G_j)|\to\infty$, istnieje jedna
liczba całkowita $d\geq1$, dla której
wszystkie $G_j$ są skończone, nieskierowane
i $d$-regularne,
oraz istnieje stała $\varepsilon>0$
taka, że $\Phi(G_j)\geq\varepsilon$
dla każdego $j$.
\end{definicja}

Zastrzeżenie o \emph{wspólnym} stopniu jest istotne:
grafy pełne mają bardzo dobrą ekspansję,
ale stopień $K_n$ rośnie wraz z $n$.
Najpierw badamy grafy regularne, następnie przejdziemy
do ważonej objętości i laplasjanu znormalizowanego
dla grafów nieregularnych.

\section{Laplasjan i wariacyjna charakterystyka luki}
\label{sec:tg-expanders-spectrum}

Niech odtąd $G=(V,E)$ będzie grafem prostym,
$d$-regularnym, o $n\geq2$ wierzchołkach, $d\geq1$.
Nie zakładamy na razie spójności.
Niech $A$ oznacza jego macierz sąsiedztwa,
zdefiniowaną w rozdziale \ref{ch:tg-podstawy}.

\begin{definicja}[Laplasjan znormalizowany]
\label{def:tg-expander-laplacian}
Niech $G=(V,E)$ będzie skończonym grafem prostym
nieskierowanym, $d$-regularnym, $n=|V|\geq2$, $d\geq1$,
a $A$ jego macierzą sąsiedztwa.
Definiujemy $P=A/d$ i $\mathcal L=I-P=I-A/d$.
Oznaczmy wartości własne macierzy
rzeczywistej symetrycznej $\mathcal L$
(z uwzględnieniem krotności)
przez
$0=\lambda_1\leq\lambda_2\leq\cdots\leq\lambda_n$.
Liczbę $\lambda_2$ nazywamy
\emph{znormalizowaną luką spektralną} grafu.
Zgodnie z uwagą \ref{rem:tg-laplacian-conventions}
w rozdziale \ref{ch:tg-laplacian} mamy
$\mathcal L=d^{-1}L_G$, gdzie $L_G=BB^T$.
Nie utożsamiamy więc tej luki z drugą wartością
własną macierzy $L_G$.
\end{definicja}

\begin{lemat}[Postać energetyczna]
\label{lem:tg-expander-energy}
Niech $G=(V,E)$ będzie skończonym grafem prostym
nieskierowanym, $d$-regularnym, $|V|\geq2$, $d\geq1$.
W przestrzeni $\mathbb R^V$ używamy iloczynu
skalarnego $\langle f,g\rangle=\sum_v f(v)g(v)$.
Dla każdej funkcji $f:V\to\mathbb R$ zachodzi
\[
d\langle f,\mathcal L f\rangle
=\sum_{\{u,v\}\in E}(f(u)-f(v))^2.
\]
Wszystkie wartości własne $\mathcal L$
należą do $[0,2]$, a wymiar jądra
$\mathcal L$ jest liczbą składowych spójnych $G$.
W szczególności
$\lambda_2>0$ dokładnie wtedy, gdy $G$ jest spójny.
\end{lemat}
\begin{proof}
Rozwijając kwadraty i wykorzystując
$d$-regularność, otrzymujemy
$\sum_{\{u,v\}\in E}(f(u)-f(v))^2
=d\sum_vf(v)^2-2\sum_{\{u,v\}\in E}f(u)f(v)$,
co jest równe $d\langle f,(I-A/d)f\rangle$.
Wynik jest nieujemny.
Z nierówności $(a-b)^2\leq2(a^2+b^2)$
otrzymujemy górne oszacowanie
$\langle f,\mathcal L f\rangle\leq2\|f\|_2^2$,
a więc widmo zawiera się w $[0,2]$.
Energia znika dokładnie wtedy,
gdy $f$ jest stała na każdej krawędzi,
czyli na każdej składowej.
Ponieważ $\mathcal L$ jest symetryczna i dodatnio
półokreślona, zerowa energia jest równoważna
przynależności do jej jądra (rozwińmy funkcję w bazie
wektorów własnych). Funkcje charakterystyczne składowych
tworzą bazę tego jądra, a krotność zera równa się
jego wymiarowi.
\end{proof}

\begin{stwierdzenie}[Iloraz Rayleigha]
\label{prop:tg-expander-rayleigh}
Dla każdego skończonego grafu prostego
$d$-regularnego $G$ o $n\geq2$, $d\geq1$,
\[
\lambda_2=
\min_{\substack{f:V\to\mathbb R,\ f\not\equiv0\\
                 \sum_{v\in V}f(v)=0}}
\frac{\sum_{\{u,v\}\in E}(f(u)-f(v))^2}
{d\sum_{v\in V}f(v)^2}.
\]
\end{stwierdzenie}
\begin{proof}
Wektor stały $\boldsymbol1$ jest
wektorem własnym $\mathcal L$
z wartością własną $0$.
Z twierdzenia spektralnego
można wybrać ortonormalną
bazę wektorów własnych,
w której pierwszym jest
$\boldsymbol1/\sqrt n$.
Warunek $\sum f(v)=0$
oznacza ortogonalność do
tego wektora. Iloraz energii
i normy wektora jest wtedy
średnią ważoną wartości
$\lambda_2,\ldots,\lambda_n$
i jest co najmniej $\lambda_2$.
Równość uzyskujemy dla
odpowiedniego wektora własnego.
Korzystamy z lematu
\ref{lem:tg-expander-energy}.
\end{proof}

\section{Nierówności Cheegera dla grafów regularnych}
\label{sec:tg-expanders-cheeger}

\begin{lemat}[Nierówność współpowierzchniowa dla funkcji]
\label{lem:tg-expander-coarea}
Niech $G=(V,E)$ będzie skończonym grafem
$d$-regularnym o $d\geq1$, $|V|=n\geq2$.
Niech $g:V\to[0,\infty)$ będzie niezerową
funkcją, której nośnik ma co najwyżej $n/2$
elementów. Wówczas
\[
\sum_{\{u,v\}\in E}(g(u)-g(v))^2
\geq\frac{d\Phi(G)^2}{2}\sum_{v\in V}g(v)^2.
\]
\end{lemat}
\begin{proof}
Dla $t\geq0$ niech
$S_t=\{v:g(v)^2>t\}$.
Każdy niepusty $S_t$
zawiera co najwyżej $n/2$ wierzchołków;
z definicji $\Phi(G)$ mamy
$|\partial_ES_t|\geq d\Phi(G)|S_t|$.
Dla każdej krawędzi $\{u,v\}$ zbiór parametrów $t$,
dla których krawędź należy do $\partial_E S_t$,
jest przedziałem pomiędzy $g(u)^2$ i $g(v)^2$
(z dokładnością do końców). Jego długość wynosi
$|g(u)^2-g(v)^2|$. Analogicznie całka funkcji
charakterystycznej warunku $v\in S_t$ wynosi
$g(v)^2$. Sumując te tożsamości po skończonych
zbiorach krawędzi i wierzchołków oraz oznaczając
miarę Lebesgue'a na prostej przez $\lambda_{\mathbb R}$ (aby nie
mylić jej z wartością własną $\lambda_1$), otrzymujemy
\[
\begin{aligned}
\sum_{\{u,v\}\in E}|g(u)^2-g(v)^2|
&=\int_0^\infty|\partial_ES_t|\,\mathrm d\lambda_{\mathbb R}(t)\\
&\geq d\Phi(G)\int_0^\infty|S_t|\,\mathrm d\lambda_{\mathbb R}(t)\\
&=d\Phi(G)\sum_{v\in V}g(v)^2.
\end{aligned}
\]
Z drugiej strony nierówności
Cauchy'ego--Schwarza i
$(a+b)^2\leq2(a^2+b^2)$ dają
\[
\begin{aligned}
\sum_{\{u,v\}\in E}|g(u)^2-g(v)^2|
&\leq
\left(\sum_{\{u,v\}\in E}(g(u)-g(v))^2\right)^{1/2}\\
&\quad{}\cdot
\left(\sum_{\{u,v\}\in E}(g(u)+g(v))^2\right)^{1/2}\\
&\leq
\left(2d\sum_v g(v)^2\right)^{1/2}
\left(\sum_{\{u,v\}\in E}(g(u)-g(v))^2\right)^{1/2}.
\end{aligned}
\]
Podnosząc do kwadratu
i skracając dodatni czynnik
$d\sum_vg(v)^2$, otrzymujemy
żądaną nierówność.
\end{proof}

\begin{twierdzenie}[Nierówności Cheegera]
\label{thm:tg-expander-cheeger}
Dla każdego grafu prostego
$d$-regularnego $G$
o $n\geq2$ wierzchołkach i $d\geq1$
zachodzi
\[
\frac{\lambda_2(G)}2
\leq\Phi(G)\leq\sqrt{2\lambda_2(G)}.
\]
Równoważnie:
$\Phi(G)^2/2\leq\lambda_2(G)\leq2\Phi(G)$.
\end{twierdzenie}
\begin{proof}
Wybierzmy niepusty zbiór
$S\subset V$ o $|S|\leq n/2$
realizujący minimum $\Phi(G)$.
Dla $f=\chi_S-(|S|/n)\boldsymbol1$
mamy $\sum_vf(v)=0$ oraz
$\sum_vf(v)^2=|S|(n-|S|)/n$.
Energia $f$ równa się
$|\partial_ES|$, więc iloraz
Rayleigha z
\ref{prop:tg-expander-rayleigh} daje
\[
\lambda_2
\leq\frac{|\partial_ES|}
{d|S|(n-|S|)/n}
=\Phi(G)\frac n{n-|S|}
\leq2\Phi(G).
\]

Do przeciwnej nierówności
wybierzmy rzeczywisty
wektor własny $f\not\equiv0$
dla $\lambda_2$, prostopadły
do wektora stałego.
Niezerowa funkcja o sumie wartości zero
przyjmuje zarówno wartości dodatnie, jak i ujemne.
Wybierzmy znak $f$ tak, aby co najwyżej $n/2$
jego wartości było dodatnich
(jeden z dwóch zbiorów:
$\{f>0\}$, $\{f<0\}$,
musi mieć ten rozmiar).
Niech $g(v)=\max\{f(v),0\}$.
Funkcja $g$ jest niezerowa,
nieujemna i ma nośnik
o co najwyżej $n/2$ elementach.

Dla wszystkich liczb rzeczywistych
$a,b$ zachodzi
$(a-b)(a_+-b_+)\geq(a_+-b_+)^2$,
gdzie $x_+=\max(x,0)$: funkcja $x\mapsto x_+$
jest niemalejąca i $1$-lipschitzowska, więc różnice
mają zgodny znak, a $|a_+-b_+|\leq|a-b|$.
Polaryzacja tożsamości z lematu
\ref{lem:tg-expander-energy} daje
$d\langle g,\mathcal Lf\rangle=
\sum_{\{u,v\}\in E}(f(u)-f(v))(g(u)-g(v))$.
Ponieważ $f(v)g(v)=g(v)^2$, po pomnożeniu
równania własnego przez $g$ otrzymujemy
\[
\lambda_2d\sum_vg(v)^2
=\sum_{\{u,v\}\in E}(f(u)-f(v))(g(u)-g(v))
\geq\sum_{\{u,v\}\in E}(g(u)-g(v))^2.
\]
Lemat \ref{lem:tg-expander-coarea}
daje dolne ograniczenie
ostatniej sumy przez
$d\Phi(G)^2\sum_vg(v)^2/2$.
Dzieląc przez
$d\sum_vg(v)^2>0$,
dostajemy
$\lambda_2\geq\Phi(G)^2/2$.
\end{proof}

\begin{wniosek}[Ekspansja a luka spektralna w rodzinie]
\label{cor:tg-expander-spectral-family}
Niech $(G_j)_{j\geq1}$ będzie rodziną
grafów prostych $d$-regularnych o wspólnym
$d\geq1$, przy czym $|V(G_j)|\to\infty$.
Jest to rodzina ekspanderów w sensie
definicji \ref{def:tg-expander-family}
wtedy i tylko wtedy, gdy
$\inf_{j\geq1}\lambda_2(G_j)>0$.
\end{wniosek}
\begin{proof}
Jeżeli $\Phi(G_j)\geq\varepsilon>0$,
to nierówności Cheegera dają
$\lambda_2(G_j)\geq\varepsilon^2/2$.
Odwrotnie, jeżeli
$\lambda_2(G_j)\geq\delta>0$,
to $\Phi(G_j)\geq\delta/2$.
\end{proof}

\section{Lemat o mieszaniu krawędzi}
\label{sec:tg-expander-edge-mixing}

Ekspansja zbioru kontroluje jego przekrój z dopełnieniem.
Widmo pozwala ponadto porównać liczbę połączeń między
dwoma dowolnymi zbiorami z liczbą, jakiej oczekiwalibyśmy
przy równomiernym rozmieszczeniu krawędzi.

\begin{twierdzenie}[Lemat o mieszaniu ekspanderów]
\label{thm:tg-expander-edge-mixing}
Niech $G=(V,E)$ będzie skończonym grafem prostym
nieskierowanym, $d$-regularnym, $n=|V|\geq2$, $d\geq1$.
Niech $d=\theta_1\geq\theta_2\geq\cdots\geq\theta_n$
będą wartościami własnymi macierzy sąsiedztwa $A$,
a $\Lambda=\max_{2\leq i\leq n}|\theta_i|$.
Dla dowolnych $S,T\subseteq V$ połóżmy
$e(S,T)=\sum_{u\in S,v\in T}A_{uv}$.
Wówczas
\[
\left|e(S,T)-\frac{d|S||T|}{n}\right|
\leq\Lambda
\sqrt{|S|(1-|S|/n)|T|(1-|T|/n)}.
\]
Jeżeli $S,T$ są rozłączne, $e(S,T)$ jest liczbą
krawędzi między nimi; gdy $S=T$, każda wewnętrzna
krawędź jest liczona dwukrotnie.
\end{twierdzenie}
\begin{proof}
Piszemy $\chi_S=(|S|/n)\boldsymbol1+f$ oraz
$\chi_T=(|T|/n)\boldsymbol1+g$, gdzie
$f,g$ są prostopadłe do $\boldsymbol1$.
Ponieważ $A\boldsymbol1=d\boldsymbol1$ i $A$ jest
symetryczna, składniki mieszane znikają:
\[
e(S,T)=\langle\chi_S,A\chi_T\rangle
=\frac{d|S||T|}{n}+\langle f,Ag\rangle.
\]
Na podprzestrzeni prostopadłej do $\boldsymbol1$
norma operatorowa $A$ wynosi $\Lambda$ na mocy
twierdzenia spektralnego, również gdy wartość własna
$d$ ma krotność większą niż jeden.
Stąd $|\langle f,Ag\rangle|\leq\Lambda\|f\|_2\|g\|_2$.
Bezpośrednie rozwinięcie kwadratów daje
$\|f\|_2^2=|S|-|S|^2/n$ i analogiczny wzór dla $g$,
co kończy dowód.
\end{proof}

\begin{uwaga}[Dwa różne parametry spektralne]
Znormalizowana luka $\lambda_2=1-\theta_2/d$
kontroluje ekspansję i mieszanie leniwego spaceru.
Lemat o mieszaniu krawędzi wykorzystuje jednak
również dolny kraniec widma: $\Lambda/d=
\max\{1-\lambda_2,\lambda_n-1\}$.
Na spójnym grafie dwudzielnym wektor o wartościach
$1$ i $-1$ na dwóch częściach jest wektorem własnym
$A$ dla $-d$, więc $\Lambda=d$ mimo dodatniej
luki $\lambda_2$. Nie należy utożsamiać tych parametrów.
\end{uwaga}

\section{Leniwy spacer losowy i tempo mieszania}
\label{sec:tg-expanders-random-walk}

W przypadku grafów regularnych mieszanie oznacza
zbieżność \emph{rozkładu położenia} spaceru losowego
do rozkładu jednostajnego na wierzchołkach;
nie chodzi o losowy wybór krawędzi
w modelu $G(n,p)$.

\begin{definicja}[Leniwy spacer losowy]
\label{def:tg-expander-lazy-walk}
Niech $G=(V,E)$ będzie spójnym grafem
prostym $d$-regularnym o $n\geq2$
wierzchołkach, $d\geq1$.
Niech $P=A/d$ i
$Q=(I+P)/2$.
Łańcuch Markowa o macierzy przejścia
$Q$ nazywamy \emph{leniwym spacerem losowym}:
w każdym kroku pozostaje w miejscu
z prawdopodobieństwem $1/2$,
a z prawdopodobieństwem $1/2$
wybiera jednostajnie jednego
z $d$ sąsiadów.
\end{definicja}

Ponieważ $Q$ jest symetryczna
i sumy jej wierszy wynoszą $1$,
rozkład $\pi(v)=1/n$ jest
stacjonarny. Dla dwóch rozkładów
$\mu,\nu$ na $V$ definiujemy
odległość w wariacji całkowitej
\[
\|\mu-\nu\|_{\mathrm{TV}}
=\frac12\sum_{v\in V}|\mu(v)-\nu(v)|.
\]
Wektor rozkładu utożsamiamy
z wektorem w $\mathbb R^V$.

\begin{twierdzenie}[Mieszanie z prędkością spektralną]
\label{thm:tg-expander-mixing}
Niech $G$ będzie spójnym grafem
prostym $d$-regularnym o $n\geq2$
wierzchołkach, $d\geq1$.
Niech $Q$ będzie macierzą
leniwego spaceru, $\pi$
rozkładem jednostajnym, a
$\lambda_2$ drugą wartością
własną $\mathcal L=I-A/d$.
Dla każdego rozkładu początkowego
$\mu$ i liczby całkowitej $t\geq0$
(przyjmujemy $a^0=1$ także dla $a=0$)
zachodzi
\[
\|\mu Q^t-\pi\|_{\mathrm{TV}}
\leq\frac{\sqrt n}{2}
\left(1-\frac{\lambda_2}{2}\right)^t
\|\mu-\pi\|_2
\leq\frac{\sqrt{n-1}}2
\left(1-\frac{\lambda_2}{2}\right)^t.
\]
W szczególności dla każdego
$\varepsilon\in(0,1)$
odległość ta jest co najwyżej
$\varepsilon$, gdy
\[
t\geq\max\left\{0,\,
\frac2{\lambda_2}
\log\frac{\sqrt n}{2\varepsilon}\right\}.
\]
\end{twierdzenie}
\begin{proof}
Macierz $Q=I-\mathcal L/2$
jest symetryczna, a jej
wartości własne wynoszą
$\eta_i=1-\lambda_i/2$.
Na mocy
\ref{lem:tg-expander-energy}
wszystkie należą do $[0,1]$,
a spójność daje
$\eta_2=1-\lambda_2/2<1$.
Podprzestrzeń prostopadła
do wektora stałego jest
$Q$-niezmiennicza.
Ponieważ $\mu-\pi$
ma sumę współrzędnych zero,
rozkład spektralny daje
\[
\|(\mu-\pi)Q^t\|_2
\leq\eta_2^t\|\mu-\pi\|_2.
\]
Z nierówności Cauchy'ego--Schwarza
$\|x\|_1\leq\sqrt n\|x\|_2$
otrzymujemy pierwszą nierówność
dla wariacji całkowitej.
Następnie
$\|\mu-\pi\|_2^2
=\sum_v\mu(v)^2-1/n
\leq1-1/n$,
ponieważ $\sum_v\mu(v)^2
\leq(\sum_v\mu(v))^2=1$.
Daje to drugą nierówność.

Wreszcie
$1-\lambda_2/2
\leq\exp(-\lambda_2/2)$.
Gdy wskazany dolny warunek
na $t$ zachodzi,
$\sqrt n\exp(-t\lambda_2/2)/2
\leq\varepsilon$,
a wcześniejsze oszacowania
kończą dowód.
\end{proof}

\begin{uwaga}[Dlaczego spacer jest leniwy]
\label{rem:tg-expander-lazy-needed}
Zwykły spacer na grafie
dwudzielnym może przechodzić
na przemian między dwiema
częściami. Na przykład
na $K_2$, jeśli spacer zaczyna się w ustalonym
wierzchołku, rozkład położenia w każdej chwili
jest skupiony w jednym wierzchołku
i nie zbiega do
rozkładu jednostajnego.
Macierz $Q=(I+P)/2$
eliminuje tę przeszkodę:
jej wartości własne
są nieujemne,
a dla grafu spójnego
poza stałą wszystkie
są ściśle mniejsze od $1$.
\end{uwaga}

\section{Grafy nieregularne: objętość, przewodność i mieszanie}
\label{sec:tg-expander-irregular}

Niech $G=(V,E)$ będzie skończonym spójnym grafem
prostym nieskierowanym, $n=|V|\geq2$.
Każdy stopień $d_v=d_G(v)$ jest dodatni.
Niech $D=\operatorname{diag}(d_v)$ oraz
$W=\sum_{v\in V}d_v=2|E|$.

\begin{definicja}[Objętość i przewodność]
\label{def:tg-expander-volume}
Niech $G=(V,E)$ będzie skończonym spójnym grafem
prostym nieskierowanym, $|V|\geq2$, $d_v=d_G(v)$,
$D=\operatorname{diag}(d_v)$, $W=\sum_v d_v$.
Dla $S\subseteq V$ definiujemy
$\operatorname{vol}(S)=\sum_{v\in S}d_v$ oraz
\[
\Phi_{\mathrm{vol}}(G)=
\min_{0<\operatorname{vol}(S)\leq W/2}
\frac{|\partial_E S|}{\operatorname{vol}(S)}.
\]
Definiujemy również symetryczny laplasjan
znormalizowany $\mathcal L_{\mathrm{sym}}=
I-D^{-1/2}AD^{-1/2}$.
W minimum rozważamy podzbiory $S\subseteq V$.
Rodzina jest niepusta, ponieważ przynajmniej jeden
wierzchołek ma stopień co najwyżej $W/2$.
Dla grafów regularnych obie definicje sprowadzają się
do poprzednich.
\end{definicja}

\begin{lemat}[Ważony iloraz Rayleigha]
\label{lem:tg-expander-weighted-rayleigh}
Niech $\lambda_2^{\mathrm{sym}}$ będzie drugą najmniejszą
wartością własną $\mathcal L_{\mathrm{sym}}$ dla grafu
jak w definicji \ref{def:tg-expander-volume}. Wówczas
\[
\lambda_2^{\mathrm{sym}}=
\min_{\substack{f:V\to\mathbb R,\ f\not\equiv0\\\sum_v d_v f(v)=0}}
\frac{\sum_{\{u,v\}\in E}(f(u)-f(v))^2}
{\sum_v d_v f(v)^2}.
\]
Ponadto $0<\lambda_2^{\mathrm{sym}}\leq2$.
\end{lemat}
\begin{proof}
Podstawiamy $x=D^{1/2}f$. Bezpośrednie rozwinięcie
daje $\langle x,\mathcal L_{\mathrm{sym}}x\rangle
=\sum_{\{u,v\}\in E}(f(u)-f(v))^2$ oraz
$\|x\|_2^2=\sum_v d_v f(v)^2$.
Energia jest nieujemna i znika dokładnie dla funkcji
stałych dzięki spójności. Zatem jądro jest rozpięte
przez $D^{1/2}\boldsymbol1$.
Oszacowanie $(f(u)-f(v))^2\leq2(f(u)^2+f(v)^2)$
daje górną granicę $2$ dla widma.
Ortogonalność $x$ do jądra jest równoważna
warunkowi $\sum_vd_vf(v)=0$.
Twierdzenie spektralne daje tezę.
\end{proof}

\begin{twierdzenie}[Ważone nierówności Cheegera]
\label{thm:tg-expander-weighted-cheeger}
Dla każdego skończonego spójnego grafu prostego
nieskierowanego o co najmniej dwóch wierzchołkach
zachodzi
\[
\frac{\Phi_{\mathrm{vol}}(G)^2}{2}
\leq\lambda_2^{\mathrm{sym}}\leq2\Phi_{\mathrm{vol}}(G).
\]
\end{twierdzenie}
\begin{proof}
Niech $S$ realizuje minimum przewodności,
$a=\operatorname{vol}(S)\leq W/2$.
Dla $f=\chi_S-(a/W)\boldsymbol1$ ważona średnia
wynosi zero, mianownik ilorazu Rayleigha jest
równy $a(1-a/W)$, a licznik $|\partial_ES|$.
Stąd $\lambda_2^{\mathrm{sym}}\leq
\Phi_{\mathrm{vol}}/(1-a/W)\leq2\Phi_{\mathrm{vol}}$.

Dla przeciwnego oszacowania wybieramy niezerowy
wektor własny $x$ dla $\lambda_2^{\mathrm{sym}}$,
prostopadły do $D^{1/2}\boldsymbol1$,
i kładziemy $f=D^{-1/2}x$. Wtedy
$(D-A)f=\lambda_2^{\mathrm{sym}}Df$,
a $\sum_vd_vf(v)=0$. Funkcja $f$ przyjmuje
wartości obu znaków. Po ewentualnej zmianie znaku
zbiór $\{f>0\}$ ma objętość co najwyżej $W/2$.
Niech $g=\max\{f,0\}$.
Z tożsamości
$\langle g,(D-A)f\rangle=
\sum_{\{u,v\}\in E}(g(u)-g(v))(f(u)-f(v))$,
wynikającej z rozwinięcia obu stron, i z równości
$f(v)g(v)=g(v)^2$ oraz nierówności dla części dodatnich
z dowodu twierdzenia \ref{thm:tg-expander-cheeger}
otrzymujemy
\[
\lambda_2^{\mathrm{sym}}\sum_vd_vg(v)^2
\geq\sum_{\{u,v\}\in E}(g(u)-g(v))^2.
\]
Dla $S_t=\{v:g(v)^2>t\}$ mamy
$|\partial_ES_t|\geq\Phi_{\mathrm{vol}}\operatorname{vol}(S_t)$,
także gdy $S_t$ jest pusty. Całkując względem miary Lebesgue'a $\lambda_{\mathbb R}$
na $[0,\infty)$, jak w dowodzie lematu
\ref{lem:tg-expander-coarea}, dostajemy
\[
\begin{aligned}
\sum_{\{u,v\}\in E}|g(u)^2-g(v)^2|
&=\int_0^\infty|\partial_ES_t|\,\mathrm d\lambda_{\mathbb R}(t)\\
&\geq\Phi_{\mathrm{vol}}\int_0^\infty
\sum_v d_v\chi_{S_t}(v)\,\mathrm d\lambda_{\mathbb R}(t)\\
&=\Phi_{\mathrm{vol}}\sum_vd_vg(v)^2.
\end{aligned}
\]
Z drugiej strony nierówność Cauchy'ego--Schwarza
oraz $(a+b)^2\leq2(a^2+b^2)$ dają oszacowanie
tej sumy z góry przez
\[
\left(\sum_{\{u,v\}\in E}(g(u)-g(v))^2\right)^{1/2}
\left(2\sum_v d_vg(v)^2\right)^{1/2}.
\]
Po podniesieniu do kwadratu i podzieleniu przez
$2\sum_vd_vg(v)^2>0$ uzyskujemy dolną granicę
energii $g$ równą
$\tfrac12\Phi_{\mathrm{vol}}^2\sum_vd_vg(v)^2$.
Łącząc obie nierówności dla energii i dzieląc
przez dodatnią normę ważoną, otrzymujemy tezę.
\end{proof}

\begin{twierdzenie}[Mieszanie spaceru na grafie nieregularnym]
\label{thm:tg-expander-weighted-mixing}
Niech $G=(V,E)$ będzie skończonym spójnym grafem
prostym nieskierowanym o co najmniej dwóch
wierzchołkach. Przy oznaczeniach z definicji
\ref{def:tg-expander-volume} połóżmy
$P=D^{-1}A$, $Q=(I+P)/2$ i $\pi(v)=d_v/W$.
Rozkład $\pi$ jest jedynym rozkładem stacjonarnym
dla $Q$. Dla każdego rozkładu $\mu$ na $V$ i
każdej liczby całkowitej $t\geq0$ zachodzi
\[
\|\mu Q^t-\pi\|_{\mathrm{TV}}
\leq\frac12(1-\lambda_2^{\mathrm{sym}}/2)^t
\left(\sum_v\frac{(\mu(v)-\pi(v))^2}{\pi(v)}\right)^{1/2}.
\]
Przyjmujemy $0^0=1$ dla $t=0$.
W szczególności ostatni czynnik nie przekracza
$\sqrt{1/\pi_{\min}-1}$, gdzie
$\pi_{\min}=\min_{v\in V}\pi(v)>0$.
\end{twierdzenie}
\begin{proof}
Zachodzi szczegółowa równowaga
$\pi(u)Q_{uv}=\pi(v)Q_{vu}$, gdyż dla $u\ne v$
obie strony są równe $A_{uv}/(2W)$, a dla $u=v$
równość jest natychmiastowa. Sumując po $u$,
otrzymujemy $\pi Q=\pi$.
Dla dowolnego rozkładu $\nu$ określmy jego
wycentrowaną gęstość $h_\nu(v)=\nu(v)/\pi(v)-1$.
Rozkłady traktujemy jako wektory wierszowe,
a funkcje jako wektory kolumnowe. Z równowagi
szczegółowej i sumowania wierszy $Q$ do jedności
wynika
\[
\begin{aligned}
h_{\nu Q}(v)
&=\sum_u\frac{\nu(u)Q_{uv}}{\pi(v)}-1\\
&=\sum_u Q_{vu}\bigl(h_\nu(u)+1\bigr)-1\\
&=(Qh_\nu)(v).
\end{aligned}
\]
Na przestrzeni funkcji z normą
$\|h\|_\pi^2=\sum_v\pi(v)h(v)^2$
operator $Q$ jest sprzężony izometrycznie
przez $h\mapsto W^{-1/2}D^{1/2}h$
z symetryczną macierzą
$D^{1/2}QD^{-1/2}=
(I+D^{-1/2}AD^{-1/2})/2=I-\mathcal L_{\mathrm{sym}}/2$.
Ponieważ $\sum_v\pi(v)h_\nu(v)=0$,
obraz $h_\nu$ jest prostopadły do
$D^{1/2}\boldsymbol1$. Na tej podprzestrzeni
norma operatorowa $Q^t$ nie przekracza
$(1-\lambda_2^{\mathrm{sym}}/2)^t$.
Zatem $\|h_{\mu Q^t}\|_\pi\leq
(1-\lambda_2^{\mathrm{sym}}/2)^t\|h_\mu\|_\pi$.
Nierówność Cauchy'ego--Schwarza daje
$\sum_v\pi(v)|h(v)|\leq\|h\|_\pi$, co dowodzi
oszacowania wariacji całkowitej.
Ponadto
$\|h_\mu\|_\pi^2=\sum_v\mu(v)^2/\pi(v)-1
\leq1/\pi_{\min}-1$.
Ponieważ $\lambda_2^{\mathrm{sym}}>0$,
każdy rozkład początkowy zbiega do $\pi$;
jeśli był stacjonarny, musi więc być równy $\pi$.
\end{proof}

\begin{przyklad}[Gwiazda a wybór rozkładu stacjonarnego]
\label{ex:tg-expander-star}
Niech $G=K_{1,r}$, $r\geq2$, z centrum $c$.
Zwykły spacer ma macierz przejścia $P=D^{-1}A$.
Rozkład $\pi(v)=d_v/W$ jest stacjonarny, ponieważ
$\pi(u)P_{uv}=A_{uv}/W=\pi(v)P_{vu}$ i sumując
po $u$ otrzymujemy $(\pi P)(v)=\pi(v)$.
W gwieździe $\pi(c)=1/2$, a dla każdego liścia
$\pi(v)=1/(2r)$; rozkład nie jest jednostajny.
Dopuszczalny zbiór w definicji przewodności składa
się z samych liści albo z samego centrum, ponieważ
centrum ma objętość $W/2=r$. W obu przypadkach
$|\partial_ES|=\operatorname{vol}(S)$, zatem
$\Phi_{\mathrm{vol}}(G)=1$.
\end{przyklad}

\section{Średnica i wzrost kul}
\label{sec:tg-expanders-diameter}

\begin{twierdzenie}[Średnica a ekspansja wierzchołkowa]
\label{thm:tg-expander-diameter}
Niech $G=(V,E)$ będzie spójnym grafem
prostym na $n\geq2$ wierzchołkach
i niech $h_V(G)>0$.
Wówczas
\[
\operatorname{diam}(G)\leq
2\left(
\left\lceil
\frac{\log n}{\log(1+h_V(G))}
\right\rceil+1\right).
\]
W szczególności, jeśli
$h_V(G_j)\geq\varepsilon>0$
w rodzinie grafów spójnych,
to średnice wynoszą $O(\log|V(G_j)|)$.
\end{twierdzenie}
\begin{proof}
Dla $v\in V$ niech $B_r(v)$
oznacza zbiór wierzchołków
o odległości co najwyżej $r$
od $v$. Gdy
$0<|B_r(v)|\leq n/2$,
z definicji ekspansji
mamy
\[
|B_{r+1}(v)|
=|B_r(v)|+|\partial_VB_r(v)|
\geq(1+h_V(G))|B_r(v)|.
\]
Niech
$R=\lceil\log n/\log(1+h_V(G))\rceil+1$.
Ponieważ $|B_0(v)|=1$, gdyby $|B_R(v)|\leq n/2$,
to wszystkie wcześniejsze kule
również miałyby rozmiar
co najwyżej $n/2$,
a iteracja nierówności dałaby
$|B_R(v)|\geq(1+h_V(G))^R>n$,
sprzeczność.
Zatem każda kula $B_R(v)$
zawiera więcej niż połowę
wierzchołków grafu.
Dla dowolnych $u,v$ kule
$B_R(u)$ i $B_R(v)$ mają
niepuste przecięcie.
Łącząc ich środki przez
wspólny wierzchołek,
otrzymujemy
$d_G(u,v)\leq2R$.
\end{proof}

\section{Przykłady: grafy pełne i cykle}
\label{sec:tg-expanders-examples}

\begin{przyklad}[Graf pełny]
\label{ex:tg-expander-complete}
Dla $n\geq2$ graf $K_n$
jest $(n-1)$-regularny.
Jeżeli $|S|=k\leq n/2$,
to $|\partial_ES|=k(n-k)$.
Stąd
\[
\Phi(K_n)=
\min_{1\leq k\leq n/2}
\frac{n-k}{n-1}
=\frac{\lceil n/2\rceil}{n-1}.
\]
Ponadto $A=J-I$,
gdzie $J$ jest macierzą
samych jedynek.
Na wektorze stałym
$A$ działa przez $n-1$,
a na podprzestrzeni
sumy współrzędnych zero
przez $-1$.
Zatem
$\lambda_2(K_n)=n/(n-1)$.
Mimo dobrej ekspansji
ciąg $(K_n)$ nie jest
rodziną ekspanderów
o ograniczonym stopniu,
gdyż jego stopień
$n-1$ nie jest stały.
\end{przyklad}

\begin{przyklad}[Cykl]
\label{ex:tg-expander-cycle}
Niech $n\geq3$.
Graf $C_n$ jest $2$-regularny.
Dla niepustego
$S\subsetneq V$, $|S|=k$,
granica $\partial_ES$
zawiera co najmniej
dwie krawędzie;
równość otrzymujemy,
gdy $S$ składa się
z $k$ kolejnych
wierzchołków cyklu.
Stąd
\[
\Phi(C_n)=\frac1{\lfloor n/2\rfloor}.
\]
Aby obliczyć widmo,
numerujemy wierzchołki
$0,\ldots,n-1$
w kolejności cyklicznej.
Wektory zespolone
$z_j(v)=\exp(2\pi ijv/n)$
są wektorami własnymi
macierzy $A$:
$(Az_j)(v)=z_j(v-1)+z_j(v+1)
=2\cos(2\pi j/n)z_j(v)$.
Wektory $z_j$ dla $j=0,\ldots,n-1$ są parami
ortogonalne w $\mathbb C^n$: dla $j\ne k$ iloczyn
skalarny jest sumą $n$ kolejnych potęg nietrywialnego
$n$-tego pierwiastka z jedności, równą zero.
Tworzą więc bazę. Wartości własne $\mathcal L$ wynoszą
$1-\cos(2\pi j/n)$, a najmniejsza dodatnia odpowiada
$j=1$ lub $j=n-1$. Otrzymujemy
\[
\lambda_2(C_n)
=1-\cos(2\pi/n).
\]
Zarówno $\Phi(C_n)$,
jak i $\lambda_2(C_n)$
dążą do zera;
cykle nie tworzą
rodziny ekspanderów.
\end{przyklad}



\section{Kiedy graf jest ekspanderem?}
\label{sec:tg-expanders-criteria}

\begin{figure}[htbp]
\centering
\begin{tikzpicture}[scale=.95]
\foreach \i in {0,...,7}{
  \pgfmathsetmacro{\ang}{90-45*\i}
  \node[tgVertex] (v\i) at (\ang:1.72) {\scriptsize$\i$};
}
\foreach \i in {0,...,7}{
  \pgfmathtruncatemacro{\j}{mod(\i+1,8)}
  \draw[tgEdge] (v\i)--(v\j);
}
\draw[tgAccent,line width=1.35pt]
  (v2)--(v3);
\draw[tgAccent,line width=1.35pt]
  (v7)--(v0);
\node[font=\small] at (0,-2.23)
  {$S=\{0,1,2\},\quad|\partial_ES|=2$};
\end{tikzpicture}
\caption{W cyklu $C_8$ zbiór trzech kolejnych
wierzchołków ma dwie krawędzie graniczne,
zaznaczone grubszą linią.
Dla większych cykli stosunek
$|\partial_ES|/(2|S|)$ może być coraz mniejszy.}
\label{fig:tg-expander-cycle-cut}
\end{figure}


Dla rodzin grafów $d$-regularnych
o wspólnym, stałym $d$
można stosować następujące
równoważne kryteria:
jednostajnie dodatnia
znormalizowana ekspansja
krawędziowa $\Phi$,
jednostajnie dodatnia
ekspansja wierzchołkowa $h_V$
oraz jednostajnie dodatnia
luka $\lambda_2$.
Równoważność pierwszego
i trzeciego kryterium
wynika z wniosku
\ref{cor:tg-expander-spectral-family},
a pierwszego i drugiego
ze stwierdzenia
\ref{prop:tg-expander-edge-vertex}.

W takim przypadku
twierdzenie
\ref{thm:tg-expander-mixing}
daje czas mieszania
leniwego spaceru
rzędu $O(\log n)$
dla ustalonej dokładności
w wariacji całkowitej,
a twierdzenie
\ref{thm:tg-expander-diameter}
daje średnicę $O(\log n)$.
Same małe średnice
\emph{nie stanowią}
kryterium ekspandera:
nie kontrolują
liczby krawędzi
oddzielających
dowolny mały zbiór
od jego dopełnienia.

\begin{przyklad}[Mała średnica bez jednostajnej ekspansji]
\label{ex:tg-expander-small-diameter}
Dla każdej liczby całkowitej $h\geq2$ weźmy dwie rozłączne
kopie pełnego drzewa binarnego wysokości $h$ (korzeń na
poziomie $0$, liście na poziomie $h$). W każdej kopii połączmy
jej $2^h$ liści w jeden cykl, a następnie połączmy oba
korzenie jedną krawędzią. Otrzymany graf $G_h$ jest prosty:
cykle mają co najmniej cztery wierzchołki, a żadna krawędź
drzewa nie łączy dwóch liści. Jest spójny i $3$-regularny:
wierzchołki wewnętrzne poza korzeniami mają po trzech
sąsiadów w drzewie, każdy liść ma rodzica i dwóch sąsiadów
na cyklu, a każdy korzeń ma dwoje dzieci i drugi korzeń
jako sąsiada.

Każda kopia ma $m_h=\sum_{j=0}^h2^j=2^{h+1}-1$
wierzchołków, zatem $n_h=|V(G_h)|=2m_h$. Dowolne
dwa wierzchołki w tej samej kopii można połączyć
przez jej korzeń drogą o długości co najwyżej $2h$
(po usunięciu ewentualnych powtórzeń); dla dwóch
różnych kopii przechodzimy dodatkowo przez krawędź
między korzeniami. Stąd
\[
\operatorname{diam}(G_h)\leq 2h+1=O(\log n_h).
\]
Niech $S$ będzie zbiorem wierzchołków jednej kopii.
Mamy $|S|=n_h/2$ oraz $|\partial_E S|=1$,
więc z definicji \ref{def:tg-expander-constants}
\[
0<\Phi(G_h)\leq\frac{1}{3m_h}
=\frac{2}{3n_h}\longrightarrow0.
\]
Z nierówności Cheegera \ref{thm:tg-expander-cheeger}
wynika też $\lambda_2(G_h)\leq4/(3n_h)\to0$.
Jest to zatem rodzina spójnych grafów o wspólnym
stopniu $3$ i logarytmicznej średnicy, która nie jest
rodziną ekspanderów.
\end{przyklad}

Zagadnienia jawnych konstrukcji ekspanderów,
optymalnych ograniczeń spektralnych i zastosowań
wymagają osobnego opracowania.
